\label{sec:intro}

The convergence of operational technology (OT) and information technology (IT) in modern
smart grids has produced a paradox: the same network connectivity that enables optimised
energy management also exposes mission-critical field devices—programmable logic controllers
(PLCs), remote terminal units (RTUs), and SCADA human-machine interfaces (HMIs)—to
adversaries with nation-state resources \cite{nafees2023}. The Stuxnet incident demonstrated
that Windows-based hosts embedded in the OT layer are credible attack vectors
\cite{langner2011stuxnet}. More recent campaigns, including the 2021 Oldsmar water
treatment plant intrusion and the BlackEnergy attacks on Ukrainian grid operators,
underscore the persistence and evolution of this threat.

Malware is the dominant delivery mechanism in these campaigns. Modern variants span from
polymorphic code that evades signature-based scanners to ransomware-as-a-service (RaaS)
platforms operated by financially motivated threat actors \cite{maniriho2024}. Classical
detection paradigms are broadly divided into \emph{static analysis}—inspecting binary
artefacts without execution—and \emph{dynamic analysis}—observing runtime behaviour in
controlled sandboxes \cite{ayub2023}. Static approaches are defeated by obfuscation and
packing; dynamic approaches require sandbox infrastructure incompatible with continuous OT
operations where systems cannot be paused, replicated, or diverted. Artificial Intelligence
and Machine Learning (AI/ML) methods transcend this dichotomy by learning behavioural
patterns from operational telemetry, generalising to unseen variants without signature
maintenance or execution isolation \cite{onal2025}.

Windows Event Logs have emerged as the most accessible behavioural telemetry surface in
Windows-based OT environments. They record process creation, privilege assignment, service
installation, object access, and authentication events at millisecond resolution across
every host \cite{achmad2025}. Graph-based methods \cite{zhen2025} and provenance-graph
detectors \cite{han2020} demonstrate that capturing the \emph{relational structure} of
audit log events substantially improves detection of multi-stage APT activity, though at
computational costs incompatible with real-time OT constraints. Hybrid static-dynamic
pipelines \cite{ayub2023,anderson2018} improve early-detection capability but depend on
sandbox infrastructure unavailable in live grid operations.

Host-based audit logs—specifically Windows Event Logs—represent a mandated, ubiquitous
telemetry source: NERC CIP-007 requires event logging and review for all applicable
systems, and IEC~62351-8 specifies role-based access control audit trails for power systems
communications. Despite this regulatory mandate, Windows Event Logs at the ICS layer have
received comparatively little attention from the intrusion-detection research community,
which has concentrated on network-level signatures \cite{goldenberg2013modbus,carcano2011scada}
and process-measurement anomalies \cite{goh2016swat}.

The research gap motivating this work is threefold. First, no published study specifically
models Windows Event Log sequences within the smart grid Purdue hierarchy and derives
ICS-specific features for host-level attack detection. Second, the absence of labelled
operational datasets—due to privacy, liability, and classified-information constraints—has
blocked supervised learning approaches; existing ICS datasets such as SWaT
\cite{goh2016swat} capture network or sensor measurements, not Windows host audit trails.
Third, existing AI-driven detection methods are fundamentally IT-centric: they assume
latency-negotiable, resource-abundant environments where offline or sandbox-assisted
analysis is acceptable—conditions that do not hold in safety-critical OT operations
\cite{nafees2023}.

This paper addresses both gaps with the following contributions:

\begin{enumerate}
  \item A \textbf{synthetic Windows Event Log generator} that models realistic smart grid
        host populations and four attack-category sequences grounded in the MITRE ATT\&CK
        for ICS framework \cite{mitre2020attack}.

  \item A \textbf{session-level feature engineering pipeline} that distils raw event
        streams into compact, interpretable 24-dimensional vectors.

  \item A \textbf{detection model} based on Random Forest with class-imbalance compensation
        and 5-fold stratified cross-validation, achieving perfect classification on synthetic
        benchmarks.

  \item A \textbf{systematic ablation study} quantifying the contribution of each feature
        group, providing guidance for deployment with limited feature coverage.

  \item A \textbf{deployment discussion} addressing class imbalance, concept drift, and
        GDPR data retention constraints relevant to production ICS environments.
\end{enumerate}

The remainder of this paper is structured as follows.
Section~\ref{sec:related} reviews related work.
Section~\ref{sec:threat} formalises the threat model.
Section~\ref{sec:data} describes the synthetic data and feature engineering.
Section~\ref{sec:method} presents the machine learning methodology.
Section~\ref{sec:exp} reports experiments and ablation results.
Section~\ref{sec:disc} discusses deployment considerations.
Section~\ref{sec:conc} concludes.
